\documentclass[11pt,a4paper]{article}
\usepackage{fontspec}
\setmainfont{DejaVu Sans}[Scale=0.90]
\usepackage{amsmath,amssymb}
\usepackage[margin=2.1cm]{geometry}
\usepackage{booktabs,array,tabularx}
\usepackage{enumitem}
\usepackage{titlesec}
\usepackage{parskip}
\usepackage{xcolor}
\titlelabel{\thetitle.\quad}
\titleformat{\section}{\normalfont\large\bfseries}{\thesection.}{0.6em}{}
\titleformat{\subsection}{\normalfont\normalsize\bfseries}{\thesubsection}{0.6em}{}
\titlespacing*{\section}{0pt}{15pt}{5pt}
\titlespacing*{\subsection}{0pt}{11pt}{4pt}
\setlist[itemize]{leftmargin=1.4em,itemsep=2pt,topsep=3pt}
\setlist[enumerate]{leftmargin=1.8em,itemsep=4pt,topsep=4pt}

\newcommand{\rep}[1]{\par\smallskip\noindent\textcolor{teal}{\textbf{Résultat attendu.}}\; #1\par\smallskip}

% Preserve fonts/margins; permit a final line-breaking pass for long prose.
\setlength{\emergencystretch}{2em}
\begin{document}

\begin{center}
{\Large\bfseries Exercices de partie I}\\[5pt]
{\normalsize De l'équation de Boltzmann à l'équation de Guyer--Krumhansl}\\[2pt]
{\small Version corrigée --- 15 septembre 2026}
\end{center}

\vspace{4pt}

Chaque exercice énonce une étape de la chaîne et donne le résultat à obtenir. Le calcul est laissé de
côté. Les notes 08 et 09 contiennent les définitions et le fil conducteur ; elles ne contiennent pas
les calculs intermédiaires.

L'ordre est celui de la chaîne : rien n'y est utilisé avant d'avoir été établi.

\section{Le réseau et sa dispersion}

\subsection*{Exercice 1 --- Relation de dispersion de la chaîne monoatomique}

Établir la relation de dispersion d'une chaîne d'atomes identiques de masse $M$, espacés de $a$,
reliés par des ressorts de raideur $C$, en partant de la loi de Newton appliquée au $n$-ième atome
et en cherchant des solutions en ondes planes.

\rep{$\omega(k) = 2\sqrt{C/M}\;\bigl|\sin(ka/2)\bigr|$}

\subsection*{Exercice 2 --- Limite continue et vitesse du son}

Développer le résultat précédent pour $ka \ll 1$ et identifier la vitesse du son.

\rep{$\omega \simeq vk$ avec $v = a\sqrt{C/M}$.}

\subsection*{Exercice 3 --- Vitesse de groupe au bord de zone}

Calculer $v_g = d\omega/dk$ et montrer qu'elle s'annule en $k = \pi/a$. Décrire le mouvement des
atomes voisins dans ce mode et expliquer pourquoi il ne transporte pas d'énergie.

\rep{$v_g = a\sqrt{C/M}\,\cos(ka/2)$, nulle en bord de zone. Les atomes voisins y vibrent en
opposition de phase : onde stationnaire.}

\subsection*{Exercice 4 --- Périodicité et zone de Brillouin}

Montrer que deux valeurs de $k$ différant de $2\pi/a$ décrivent le même mouvement des atomes.
Expliquer en quoi cette propriété rend possible l'existence de processus umklapp.

\section{Statistique et moments}

\subsection*{Exercice 5 --- Limites de la distribution de Bose--Einstein}

Développer $n^0 = \bigl[\exp(\hbar\omega/k_BT) - 1\bigr]^{-1}$ dans les deux limites
$\hbar\omega \ll k_BT$ et $\hbar\omega \gg k_BT$. Interpréter chacune.

\rep{$n^0 \simeq k_BT/\hbar\omega$, d'où $\hbar\omega\,n^0 \simeq k_BT$ : équipartition. Puis
$n^0 \simeq e^{-\hbar\omega/k_BT}$ : modes gelés.}

\subsection*{Exercice 6 --- Le flux est nul à l'équilibre}

Montrer que $\mathbf{q} = \int \frac{d^3k}{(2\pi)^3}\,\hbar\omega\,\mathbf{v}_g\,f$ s'annule
identiquement lorsque $f$ ne dépend que de $\omega$. Utiliser les parités respectives de $\omega$ et
de $\mathbf{v}_g$ en $\mathbf{k}$.

\rep{$\omega$ paire, $\mathbf{v}_g$ impaire, intégrande impaire, intégrale nulle. \textbf{Le flux
mesure exactement l'écart à l'équilibre.}}

\subsection*{Exercice 7 --- Proportionnalité du flux et du quasi-moment}

Pour une dispersion linéaire isotrope $\omega = vk$, montrer que $\mathbf{q} = v^2\,\mathbf{P}$, où
$\mathbf{P} = \int \frac{d^3k}{(2\pi)^3}\,\hbar\mathbf{k}\,f$.

En déduire pourquoi un processus conservant le quasi-moment ne crée aucune résistance thermique.

\rep{$\mathbf{v}_g = v\hat{\mathbf{k}}$ et $\hbar\omega = \hbar v k$, d'où
$\hbar\omega\,\mathbf{v}_g = v^2\,\hbar\mathbf{k}$. Cette identité est \textbf{la clé de toute la
suite} : elle fait disparaître les processus normaux du bilan de flux.}

\section{Hiérarchie de moments}

\subsection*{Exercice 8 --- Moment d'ordre zéro}

Multiplier l'équation de transport par $\hbar\omega_\mathbf{k}$ et intégrer sur les vecteurs d'onde.
Justifier l'annulation du terme de collision.

\rep{$\partial u/\partial t + \nabla\cdot\mathbf{q} = 0$, sans aucune hypothèse supplémentaire. Le
terme de collision s'annule parce que les deux familles de processus conservent l'énergie.}

\subsection*{Exercice 9 --- Moment d'ordre un}

Multiplier cette fois par $\hbar\omega_\mathbf{k}\,\mathbf{v}_g$ et intégrer. Identifier le moment
d'ordre deux qui apparaît. Montrer, en s'appuyant sur l'exercice~7, que seuls les processus
résistifs subsistent au second membre.

\rep{$\partial\mathbf{q}/\partial t + \nabla\cdot\mathbf{Q} = -\mathbf{q}/\tau_R$, avec
$\mathbf{Q} = \int \frac{d^3k}{(2\pi)^3}\,\hbar\omega\,\mathbf{v}_g \otimes \mathbf{v}_g\,f$.}

\subsection*{Exercice 10 --- Le problème de fermeture}

Écrire l'équation du moment d'ordre deux et constater qu'elle appelle le moment d'ordre trois.
Énoncer en une phrase pourquoi aucun niveau ne se referme de lui-même, et ce que doit fournir une
fermeture.

\subsection*{Exercice 11 --- Retrouver Fourier}

Dans l'équation du moment d'ordre un, supposer le régime stationnaire et la partie déviatorique de
$\mathbf{Q}$ négligeable. Montrer qu'on retrouve la loi de Fourier et identifier la conductivité.

\rep{$\mathbf{q} = -\lambda\nabla T$ avec $\lambda = \tfrac{1}{3}\,C\,v^2\,\tau_R$. La conductivité
est proportionnelle au temps résistif : plus les umklapp sont rares, plus le cristal conduit.}

\section{Vers Guyer--Krumhansl}

\subsection*{Exercice 12 --- Réduction unidimensionnelle}

Montrer qu'en une dimension, avec $\mathbf{q} = q(z)\,\hat{\mathbf{z}}$, les deux termes non locaux
se combinent :
\[ \nabla^2\mathbf{q} + 2\nabla(\nabla\cdot\mathbf{q}) \;\longrightarrow\;
3\,\frac{\partial^2 q}{\partial z^2} . \]

\subsection*{Exercice 13 --- Conductivité effective}

Transformer l'équation unidimensionnelle en temps, puis éliminer $\partial^2 q/\partial z^2$ à
l'aide du bilan d'énergie. En déduire que le flux s'écrit $\phi = -\lambda_{\mathrm{eff}}(p)\,\theta'$
et donner $\lambda_{\mathrm{eff}}$.

\rep{$\lambda_{\mathrm{eff}}(p) = \dfrac{\lambda + 3\ell^2 p\,\rho c}{1 + \tau_R p}$. Ce n'est pas un
simple décalage de $p$ : la structure de l'équation a changé.}

\subsection*{Exercice 14 --- Les deux combinaisons}

Poser $\xi_1 = e/\sqrt{a}$, $b = \sqrt{\lambda\rho c}$ et $\sigma^2 = p\,\rho c/\lambda_{\mathrm{eff}}$.
Montrer que
\[ \sigma e = \xi_1\sqrt{\frac{p(1+\tau_R p)}{1+\tau_\ell p}} , \qquad
\lambda_{\mathrm{eff}}\,\sigma = b\sqrt{\frac{p(1+\tau_\ell p)}{1+\tau_R p}} , \]
et identifier $\tau_\ell$.

\rep{$\tau_\ell = 3\ell^2/a$, qui est \textbf{un temps}. Les deux temps échangent leurs rôles entre
les deux expressions.}

\subsection*{Exercice 15 --- La condition de cécité}

Poser $\tau_R = \tau_\ell$ dans le résultat précédent. Conclure.

\rep{Les racines se simplifient, $\sigma e = \xi_1\sqrt{p}$ et
$\lambda_{\mathrm{eff}}\sigma = b\sqrt{p}$ : la réponse est exactement celle d'un milieu de Fourier,
quelle que soit la valeur commune des deux temps, dans le problème 1D homogène sans source intérieure ni perturbation initiale. Cette propriété est la résonance de Fourier (Kovács, 2018).}

\subsection*{Exercice 16 --- Écriture microscopique de la cécité}

À partir de $\ell^2 = v^2\tau_N\tau_R/5$ et $a = v^2\tau_R/3$, exprimer $\tau_\ell$ en fonction de
$\tau_N$ seul, puis traduire la condition de cécité.

\rep{$\tau_\ell = \tfrac{9}{5}\tau_N$, donc la condition s'écrit $\tau_R = 1{,}8\,\tau_N$ : un
rapport formel de temps de collision. À ce rapport, $\tau_N\ll\tau_R$ n'est pas satisfait ; l'extrapolation ne démontre pas une fenêtre hydrodynamique.}

\section{Admissibilité}

\subsection*{Exercice 17 --- Vitesse de propagation de Cattaneo}

Écrire le système de Cattaneo en variables $(u, q)$ sous la forme
$\partial_t U + A\,\partial_z U = 0$, calculer les valeurs propres de $A$, et conclure sur le
caractère hyperbolique.

\rep{$\pm\sqrt{a/\tau_R}$, réelles et distinctes : strictement hyperbolique. En modèle de Debye,
cette vitesse vaut $v/\sqrt{3}$.}

\subsection*{Exercice 18 --- Guyer--Krumhansl perd la vitesse finie}

À partir de la relation de dispersion $\sigma^2 = p(1+\tau_R p)/[a(1+\tau_\ell p)]$, déterminer le
comportement de $k$ à haute fréquence dans les trois cas : $\tau_R = \tau_\ell = 0$ ;
$\tau_\ell = 0$ ; les deux non nuls. Conclure sur la vitesse de phase.

\rep{$k \propto \sqrt{\omega}$, puis $k \propto \omega$, puis $k \propto \sqrt{\omega}$ à nouveau.
\textbf{Le terme non local restaure la vitesse de propagation infinie que Cattaneo avait corrigée.}}

\subsection*{Exercice 19 --- Production d'entropie de Cattaneo}

Avec l'entropie étendue $s = s_{\mathrm{éq}}(u) - \dfrac{\tau_R}{2\lambda T_0^{\,2}}\,q^2$ et le flux
d'entropie $\mathbf{J}_s = \mathbf{q}/T$, calculer $\partial_t s + \nabla\cdot\mathbf{J}_s$ en
utilisant le bilan d'énergie et la loi de Cattaneo. Développer avec $T=T_0+\varepsilon\theta$ et $q=\varepsilon j$, puis conserver les termes en $\varepsilon^2$ ; ne pas annuler les gradients en posant trop tôt $T=T_0$.

\rep{$\sigma_s = q^2/(\lambda T_0^{\,2}) \ge 0$. \textbf{Piège :} le flux d'entropie doit être pris
en $\mathbf{q}/T$ et non en $\mathbf{q}/T_0$, sans quoi subsiste un résidu en $q\,\partial_z T$ qui
ne s'annule pas.}

\subsection*{Exercice 20 --- Flux d'entropie étendu de Guyer--Krumhansl}

Reprendre le calcul précédent avec la loi de Guyer--Krumhansl, en postulant un flux d'entropie
$\mathbf{J}_s = \mathbf{q}/T + \beta\,q\,\partial_z q$. Déterminer $\beta$ en annulant le terme
croisé en $q\,\partial_z^2 q$, puis écrire la production.

\rep{$\beta = 3\ell^2/(\lambda T_0^{\,2})$, et
$\sigma_s = \bigl[q^2 + 3\ell^2(\partial_z q)^2\bigr]/(\lambda T_0^{\,2})$, somme de carrés. Sans ce
terme supplémentaire, la production n'est pas définie positive.}

\section{Contrôles}

\subsection*{Exercice 21 --- Dimensions}

Vérifier que $[\xi] = \mathrm{s}^{1/2}$, que $[V] = \mathrm{s}^{-1}$ où $V = s''/s$, et que
$\tau_\ell = 3\ell^2/a$ a bien la dimension d'un temps.

\subsection*{Exercice 22 --- Ordres de grandeur pour l'AlN}

Avec $v \approx 6\times10^3$~m/s et $\tau \approx 10^{-11}$~s, estimer le libre parcours moyen.
Comparer à l'épaisseur d'un film de 500~nm. Préciser pourquoi ce diagnostic gris illustratif ne suffit pas à valider un modèle de diffusion pour tout le spectre.

\rep{$\Lambda = v\tau \approx 60$~nm, soit environ un huitième de l'épaisseur. Les deux échelles ne
sont pas séparées : la description purement diffusive n'est pas acquise.}

\subsection*{Exercice 23 --- Fenêtre de Guyer en variables réduites}

Traduire le critère $\tau_N \ll \tau_B \ll \tau_R$ en termes de $x = \tau_R/\tau_N$ et
$y = \tau_B/\tau_N$. Décrire la région obtenue et situer la condition de cécité.

\rep{$1 \ll y \ll x$ : un coin de sommet $(1,1)$. La cécité est la verticale $x = 1{,}8$, qui passe
juste à droite du sommet géométrique. L'intervalle $1<y<x$ n'y mesure que $0{,}255$ décade et ne satisfait pas la séparation forte $1\ll y\ll x$.}

\vspace{8pt}
\noindent\small
Notes associées : \texttt{08\_\allowbreak Le\_\allowbreak gaz\_\allowbreak de\_\allowbreak phonons}, \texttt{09\_\allowbreak De\_\allowbreak Boltzmann\_\allowbreak a\_\allowbreak Guyer\_\allowbreak Krumhansl},
\texttt{11\_\allowbreak Note\_\allowbreak partie\_\allowbreak II}, \texttt{12\_\allowbreak Note\_\allowbreak partie\_\allowbreak III}. Vérifications numériques
correspondantes : \texttt{tests/\allowbreak test\_\allowbreak guyer\_\allowbreak krumhansl.py}, \texttt{tests/\allowbreak test\_\allowbreak admissibility.py}.

\end{document}
